
\documentclass{article}
\usepackage{colortbl}
\usepackage{makecell}
\usepackage{multirow}
\usepackage{supertabular}

\begin{document}

\newcounter{utterance}

\centering \large Interaction Transcript for game `eqbench', experiment `eq\_bench', episode 86 with reflexion{-}Qwen3.5{-}4B.
\vspace{24pt}

{ \footnotesize  \setcounter{utterance}{1}
\setlength{\tabcolsep}{0pt}
\begin{supertabular}{c@{$\;$}|p{.15\linewidth}@{}p{.15\linewidth}p{.15\linewidth}p{.15\linewidth}p{.15\linewidth}p{.15\linewidth}}
    \# & \multicolumn{2}{c}{Player} && \multicolumn{2}{c}{Game Master} \\
    \hline

    \theutterance \stepcounter{utterance}  
    & & & \multicolumn{4}{p{0.6\linewidth}}{
        \cellcolor[rgb]{0.9,0.9,0.9}{
            \makecell[{{p{\linewidth}}}]{
                \texttt{\tiny{[P1$\langle$GM]}}
                \\ 
\texttt{Your task is to predict the likely emotional responses of a character in this dialogue:} \\
\\ 
\texttt{Imogen: Sebastian, I wasn't expecting to see you here.} \\
\texttt{Sebastian: The feeling's mutual, Imogen.} \\
\texttt{Imogen: It's been a while.} \\
\texttt{Sebastian: Not nearly long enough, apparently.} \\
\texttt{Imogen: I see you're still as charming as ever.} \\
\texttt{Sebastian: And you're still as naive.} \\
\texttt{Imogen: That's one way to put it. I prefer to think of it as principled.} \\
\texttt{Sebastian: Principles don't pay the bills, Imogen.} \\
\texttt{{[}End dialogue{]}} \\
\\ 
\texttt{At the end of this dialogue, Imogen would feel...} \\
\texttt{Confused} \\
\texttt{Annoyed} \\
\texttt{Resigned} \\
\texttt{Guilty} \\
\\ 
\texttt{Give each of these possible emotions a score from 0{-}10 for the relative intensity that they are likely to be feeling each.} \\
\\ 
\texttt{You must output in the following format, including headings (of course, you should give your own scores), with no additional commentary:} \\
\\ 
\texttt{Confused: <score>} \\
\texttt{Annoyed: <score>} \\
\texttt{Resigned: <score>} \\
\texttt{Guilty: <score>} \\
\\ 
\\ 
\texttt{{[}End of answer{]}} \\
\\ 
\texttt{Remember: zero is a valid score, meaning they are likely not feeling that emotion. You must score at least one emotion > 0.} \\
\\ 
\texttt{Your answer:} \\
            }
        }
    }
     \\ \\

    \theutterance \stepcounter{utterance}  
    & \multicolumn{4}{p{0.6\linewidth}}{
        \cellcolor[rgb]{0.561,0.737,0.561}{
            \makecell[{{p{\linewidth}}}]{
                \texttt{\tiny{[P1$\rangle$GM]}}
                \texttt{Confused: 0} \\
\texttt{Annoyed: 5} \\
\texttt{Resigned: 0} \\
\texttt{Guilty: 0} \\
            }
        }
    }
    & & \\ \\

    \theutterance \stepcounter{utterance}  
    & & & \multicolumn{2}{p{0.3\linewidth}}{
        \cellcolor[rgb]{0.9,0.9,0.9}{
            \makecell[{{p{\linewidth}}}]{
                \texttt{\tiny{[GM$|$GM]}}
                \texttt{Confused: 0} \\
\texttt{Annoyed: 5} \\
\texttt{Resigned: 0} \\
\texttt{Guilty: 0} \\
            }
        }
    }
    & & \\ \\

    \theutterance \stepcounter{utterance}  
    & & & \multicolumn{2}{p{0.3\linewidth}}{
        \cellcolor[rgb]{0.9,0.9,0.9}{
            \makecell[{{p{\linewidth}}}]{
                \texttt{\tiny{[GM$|$GM]}}
                \texttt{Confused: 0} \\
\texttt{Annoyed: 8} \\
\texttt{Resigned: 0} \\
\texttt{Guilty: 0} \\
            }
        }
    }
    & & \\ \\

    \theutterance \stepcounter{utterance}  
    & & & \multicolumn{2}{p{0.3\linewidth}}{
        \cellcolor[rgb]{0.9,0.9,0.9}{
            \makecell[{{p{\linewidth}}}]{
                \texttt{\tiny{[GM$|$GM]}}
                \texttt{game\_result = LOSE} \\
            }
        }
    }
    & & \\ \\

    \theutterance \stepcounter{utterance}  
    & & \multicolumn{4}{p{0.6\linewidth}}{
        \cellcolor[rgb]{0.95,0.95,0.95}{
            \makecell[{{p{\linewidth}}}]{
                \texttt{\tiny{[P1$\leftrightarrow$P]}}
                \texttt{Post{-}episode Reflexion (not sent to GM):} \\
\texttt{New plan: I misinterpreted the emotional dynamic of the dialogue. Sebastian's comment "Principles don't pay the bills" is a direct, slightly condescending rebuttal to Imogen's claim of being "principled." This specific line is designed to make the listener feel that their moral stance is useless or foolish, which typically triggers feelings of **guilt** (feeling foolish or wrong) rather than annoyance or confusion. I underestimated the weight of this specific insult. Next time, I will analyze the subtext of the final exchange more deeply, recognizing that calling someone's core identity ("principled") a financial liability is a move intended to induce guilt or shame, and I will assign a high score to Guilty while keeping Annoyance lower.} \\
            }
        }
    }
    & \\ \\

\end{supertabular}
}

\end{document}
